% Chapter 11: Proto-Cohesion and the Diagnostic Vector
% Part III: Multi-Formation and Temporal Theory

\chapter{원시-응집과 진단 벡터}
\label{ch:diagnostic-vector}

\begin{quote}
\textit{``형성이 '충분히 좋은가'는 하나의 수치가 아닌 네 차원의 벡터로 판정된다.''}
\end{quote}

\vspace{0.5em}
이론의 목표는 응집적 형성의 출현과 안정성을 설명하는 것이다. 그러나 ``형성이 잘 형성되었는가?''라는 질문에는 여러 독립적인 차원이 있다. 닫힘을 잘 만족하나 외부와 구별이 약할 수 있고, 내부 구조는 좋으나 시간적으로 지속하지 못할 수 있다.

SCC는 이러한 다차원적 평가를 \textbf{원시-응집 진단 벡터(proto-cohesion diagnostic vector)} $\mathbf{d} \in [0,1]^4$로 형식화한다. 이 장에서는 네 성분 각각의 정의, 계산, 이론적 속성을 상세히 다루고, 이들이 어떤 예측을 생성하며, 그 예측들이 어떻게 검증되었는지를 살펴본다.


\section{진단 벡터 $\mathbf{d} \in [0,1]^4$}
\label{sec:diagnostic-vector}

\subsection{정의}

형성 $u_t$에 대한 원시-응집 진단 벡터는 다음과 같다:
\begin{equation}
    \mathbf{d}(u_t) = \Big(\, \mathsf{Bind}(u_t),\;\; \mathsf{Sep}(u_t),\;\; \mathsf{Inside}(u_t),\;\; \mathsf{Persist}(u_t) \,\Big) \in [0,1]^4
    \label{eq:diagnostic-vector}
\end{equation}

각 성분은 형성의 구조적 품질을 독립적인 차원에서 평가한다:
\begin{itemize}
    \item $\mathsf{Bind}$: 닫힘 하에서의 자기-지지 --- 형성이 얼마나 관계적으로 완결적인가
    \item $\mathsf{Sep}$: 외부와의 구별 --- 내부와 외부가 얼마나 비대칭적인가
    \item $\mathsf{Inside}$: 형태학적 품질 --- 내부 구조가 얼마나 명확하고 우세한가
    \item $\mathsf{Persist}$: 시간적 계승 --- 응집적 핵심이 시간을 가로질러 얼마나 보존되는가
\end{itemize}

\subsection{연속적 평가 vs. 이진 판정}

진단 벡터의 핵심적 설계 선택은 \textbf{연속적 값}이라는 점이다. 전통적으로 ``$X$가 객체인가?''는 이진 판정(예/아니오)이다. 그러나 SCC에서 ``$u_t$가 형성인가?''는 연속적 평가이다:

\begin{itemize}
    \item $\mathbf{d} = (0.95, 0.92, 0.88, 0.91)$: 모든 차원에서 강한 형성
    \item $\mathbf{d} = (0.82, 0.45, 0.71, 0.93)$: 결합과 지속은 강하나 분리가 약한 형성
    \item $\mathbf{d} = (0.31, 0.22, 0.15, 0.50)$: 약한 형성---``아직 형성이 되는 중''
\end{itemize}

이 연속적 평가가 이진 판정보다 풍부한 이유: 형성의 ``실패 모드(failure mode)''를 진단할 수 있다. ``분리가 약하다''는 진단은 ``형성이 아니다''라는 판정보다 훨씬 더 유용하다.

\subsection{왜 네 차원인가}

네 차원이 아닌 다른 선택은 어떤가?

\begin{itemize}
    \item \textbf{세 차원 (Persist 제거)}: 시간적 측면을 잃는다. 순간적으로 잘 형성되었으나 지속하지 못하는 ``잠깐의 형성''을 구별할 수 없다.
    \item \textbf{다섯 차원 ($C_t$ 기반 추가)}: v2.1에서 $C_t$가 진단적 역할로 격하되었으므로, 독립적인 다섯 번째 차원의 근거가 약하다. Sep이 이미 내외부 비대칭을 포착한다.
    \item \textbf{하나의 스칼라 (min 또는 평균)}: 차원간 트레이드오프를 숨긴다. $\min(\mathbf{d})$는 보수적 하한이고, $\mathrm{mean}(\mathbf{d})$는 약한 차원을 은폐한다.
\end{itemize}

결론: 네 차원은 에너지의 네 항(닫힘, 분리, 경계/형태학, 전달)에 정확히 대응하며, 각각이 독립적인 ``실패 축''을 포착한다. 이 대응은 우연이 아니다---에너지 설계 자체가 네 가지 구조적 요구에서 출발했기 때문이다.

\subsection{진단 벡터와 에너지의 대응 관계}

네 진단 성분과 네 에너지 항의 정확한 대응을 정리하면:

\begin{center}
\begin{tabular}{lccc}
\hline
\textbf{진단 성분} & \textbf{에너지 항} & \textbf{관계} & \textbf{강도} \\
\hline
$\mathsf{Bind}$ & $\mathcal{E}_{\mathrm{cl}}$ & $\mathsf{Bind} \geq 1 - \sqrt{\mathcal{E}_{\mathrm{cl}}/n}$ & 하한 (Cauchy-Schwarz) \\
$\mathsf{Sep}$ & $\mathcal{E}_{\mathrm{sep}}$ & $\mathsf{Sep} = 1 - \mathcal{E}_{\mathrm{sep}}/m$ & 정확한 등식 \\
$\mathsf{Inside}$ & $\mathcal{E}_{\mathrm{bd}}$ & 정성적 대응 & 간접적 \\
$\mathsf{Persist}$ & $\mathcal{E}_{\mathrm{tr}}$ & 전달 핵 기반 & 간접적 \\
\hline
\end{tabular}
\end{center}

$\mathsf{Sep}$과 $\mathcal{E}_{\mathrm{sep}}$의 관계는 정확한 등식(술어-에너지 다리)이며, 이는 이론에서 가장 아름다운 결과 중 하나이다. $\mathsf{Bind}$의 에너지 하한은 Cauchy-Schwarz에 의한 것으로, 역방향 한계(에너지 최소점에서의 상한)도 알려져 있다. $\mathsf{Inside}$와 $\mathsf{Persist}$의 대응은 간접적이나, 에너지 최소화가 네 성분 모두를 동시에 개선하는 경향이 있다.

\begin{example}[진단 벡터의 해석: 세 가지 형성 프로필]
\label{ex:three-profiles}
$20 \times 20$ 격자에서 세 가지 다른 매개변수로 형성을 최적화한 결과:

\textbf{프로필 A} ($\beta = 50$, $\lambda_{\mathrm{cl}} = 1.0$): $\mathbf{d} = (0.94, 0.91, 0.87, 0.93)$
\begin{itemize}
    \item 해석: 모든 차원에서 강한 형성. 날카로운 경계, 깊은 핵심, 강한 시간적 계승.
    \item 가장 약한 차원: $\mathsf{Inside} = 0.87$. 이유: H0 지속에서 부차적 성분이 약간 존재.
\end{itemize}

\textbf{프로필 B} ($\beta = 15$, $\lambda_{\mathrm{cl}} = 0.5$): $\mathbf{d} = (0.82, 0.45, 0.71, 0.90)$
\begin{itemize}
    \item 해석: 자기-지지는 양호하나 분리가 약함. 경계가 넓어 내외부 구별이 흐릿함.
    \item 진단: $\beta$를 높이면 $\mathsf{Sep}$이 개선될 것. $\lambda_{\mathrm{cl}}$ 증가는 $\mathsf{Bind}$를 개선하나 $\mathsf{Sep}$에는 제한적 효과.
\end{itemize}

\textbf{프로필 C} ($\beta = 5$, $\lambda_{\mathrm{cl}} = 0.1$): $\mathbf{d} = (0.31, 0.22, 0.15, 0.55)$
\begin{itemize}
    \item 해석: 전반적으로 약한 형성. 상전이 임계점($\beta_{\mathrm{crit}} \approx 8$) 이하이므로 구조화가 미약.
    \item 진단: 이 형성은 ``아직 형성이 되는 중''이며, 매개변수 조정 없이는 원시-응집 문턱을 넘기 어려움.
\end{itemize}
\end{example}


\section{Bind --- 닫힘 하의 자기-지지}
\label{sec:bind}

\subsection{정의}

\begin{equation}
    \mathsf{Bind}(u_t) = 1 - \frac{\|u_t - \mathrm{Cl}_t(u_t)\|_2}{\sqrt{n}}
    \label{eq:bind}
\end{equation}

$\mathsf{Bind}$는 응집 장 $u_t$가 닫힘 연산자 $\mathrm{Cl}_t$의 고정점에 얼마나 가까운지를 측정한다. $\mathrm{Cl}_t(u_t) = u_t$이면 $\mathsf{Bind} = 1$ (완벽한 자기-지지)이고, 잔차가 클수록 $\mathsf{Bind}$가 감소한다.

\subsection{왜 $\ell^2$ 노름이고 $\ell^\infty$가 아닌가}

직관적으로, ``최악의 지점에서 닫힘을 잘 만족하는가''를 평가하려면 $\ell^\infty$ 노름이 자연스러워 보인다. 그러나 경계 지점(boundary site)은 닫힘 잔차가 구조적으로 작지 않다: 경계 지점의 이웃은 내부(높은 $u$)와 외부(낮은 $u$)가 혼재하므로, $P_t u(x)$와 $u(x)$ 사이에 필연적 간극이 존재한다.

구체적으로, 경계 지점에서의 닫힘 잔차의 하한은 약 0.21이다. 이는 ``형성의 결함''이 아니라 경계 대역의 \emph{구조적 특성}이다. $\ell^\infty$ 노름은 이 구조적 특성을 ``실패''로 오판한다. $\ell^2$ 노름은 경계의 국소적 잔차를 전체적 평균 안에서 적절히 희석시킨다.

\subsection{T-Bind-Proj 와 T-Bind-Full}

\begin{theorem}[T-Bind-Proj]
\label{thm:t-bind-proj}
$\tau = 1/2$일 때, 응집 장 $u$를 닫힘의 고정점으로 사영하면 $\mathsf{Bind}$가 단조 증가한다:
\begin{equation}
    \mathsf{Bind}(\mathrm{Proj}_{\mathrm{FP}}(u)) \geq \mathsf{Bind}(u)
    \label{eq:bind-proj}
\end{equation}
\end{theorem}

\begin{theorem}[T-Bind-Full]
\label{thm:t-bind-full}
전체 에너지 $\mathcal{E}$의 최소점 $u^*$는 거의 최대의 $\mathsf{Bind}$를 달성한다. 에너지 한계:
\begin{equation}
    \mathsf{Bind}(u^*) \geq 1 - \sqrt{\mathcal{E}_{\mathrm{cl}}(u^*) / n}
    \label{eq:bind-energy}
\end{equation}
\end{theorem}

\subsection{일반 $\tau$에 대한 이진 질량-균형 공식}

$\tau_{\mathrm{cl}} = 1/2$가 아닌 일반적인 문턱값에 대해, $\mathsf{Bind}$의 계산은 이진 질량-균형(binary mass-balance) 분석을 통해 확장된다. 형성의 핵심은 $\sigma(a_{\mathrm{cl}}(c - \tau))$에 의해 결정되는 포화 수준을 가지며, 이 수준이 $\tau$에 의존한다. T-Bind-Proj/Full은 $\tau = 1/2$에 대해 완전히 증명되었으며, 일반 $\tau$에 대해서는 현재 카테고리 B이다.


\section{Sep --- 외부와의 구별}
\label{sec:sep}

\subsection{정의}

\begin{equation}
    \mathsf{Sep}(u_t) = \frac{\sum_{x \in X_t} u_t(x) \cdot D_t(x;\, 1-u_t)}{\sum_{x \in X_t} u_t(x)}
    \label{eq:sep}
\end{equation}

$\mathsf{Sep}$은 형성의 각 지점에서의 구별도($D_t$)를 응집 강도($u_t$)로 가중 평균한 것이다. 형성 내부의 지점일수록 높은 가중을 받으며, 이들 지점에서의 구별도가 높으면 $\mathsf{Sep}$이 높아진다.

\subsection{정확한 등식: $\mathsf{Sep} = 1 - \mathcal{E}_{\mathrm{sep}} / m$}

$\mathsf{Sep}$의 가장 주목할 만한 성질은 분리 에너지와의 정확한 등식이다:

\begin{theorem}[술어-에너지 다리, Sep 성분]
\label{thm:sep-energy-bridge}
\begin{equation}
    \mathsf{Sep}(u_t) = 1 - \frac{\mathcal{E}_{\mathrm{sep}}(u_t)}{m}
    \label{eq:sep-identity}
\end{equation}
여기서 $m = \sum_x u_t(x)$는 전체 체적이다.
\end{theorem}

이 등식은 진단적 술어(predicate)와 에너지 항 사이의 직접적인 다리를 놓는다. $\mathcal{E}_{\mathrm{sep}}$가 작을수록 $\mathsf{Sep}$이 높다는 것은, 에너지 최소화가 자연스럽게 분리를 강화함을 의미한다.

\subsection{$C_t$-가중의 실패}

초기 설계(I8 이전)에서는 공동소속 $C_t$를 가중으로 사용하는 버전이 고려되었다:
\begin{equation}
    \mathsf{Sep}_{\mathrm{old}} = \frac{\sum_x C_t(x,x) \cdot D_t(x)}{\sum_x C_t(x,x)} \quad \text{(기각됨)}
    \label{eq:sep-old}
\end{equation}
이 버전은 형성의 품질과 관계없이 약 $0.5$를 반환하는 것으로 확인되었다. 원인: $C_t(x,x)$는 외부 지점에서도 비-영(nonzero)이며, 외부 지점의 $D \approx 0$과 내부 지점의 $D \approx 1$이 평균되어 항상 중간값을 준다. $u_t$-가중은 이 문제를 해결한다---외부 지점($u \approx 0$)은 자연스럽게 무시된다.


\section{Inside --- 형태학적 품질}
\label{sec:inside}

\subsection{정의}

$\mathsf{Inside}$는 형성의 내부 구조가 얼마나 명확하고 우세한지를 측정한다:
\begin{equation}
    \mathsf{Inside}(u_t) = \mathcal{Q}_{\mathrm{morph}}(u_t) = \frac{\ell_{\max}(u_t) - c}{1 - c} \cdot \mathrm{Artic}(u_t)
    \label{eq:inside}
\end{equation}
여기서:
\begin{itemize}
    \item $\ell_{\max}$: 초수준 집합 여과(superlevel-set filtration)의 H0 지속 다이어그램에서 가장 긴 막대(bar)
    \item $c = \bar{u} = \frac{1}{n}\sum_x u(x)$: 평균 응집값
    \item $\mathrm{Artic}(u_t) = 1 - \ell_{\mathrm{second}} / \ell_{\max}$: 관절 비율(articulation ratio)
\end{itemize}

\subsection{H0 지속과 초수준 집합 여과}

$\mathsf{Inside}$의 계산은 위상적 자료 분석(Topological Data Analysis, TDA)의 도구를 사용한다. 문턱값 $\theta$를 1에서 0으로 내리면서 초수준 집합 $\{x : u(x) \geq \theta\}$의 연결 성분을 추적한다:

\begin{enumerate}
    \item $\theta = 1$: 빈 집합 (또는 $u(x) = 1$인 점만)
    \item $\theta$ 감소: 새로운 연결 성분이 ``탄생'' (높은 $u$ 지점 추가)
    \item $\theta$ 더 감소: 성분들이 ``합병'' (인접 성분이 연결됨) → 탄생한 성분이 ``소멸''
    \item $\theta = 0$: 모든 점이 포함된 하나의 연결 성분
\end{enumerate}

각 연결 성분의 ``수명''은 $\ell = \theta_{\mathrm{birth}} - \theta_{\mathrm{death}}$이다. 가장 긴 수명 $\ell_{\max}$는 ``주요 형성''의 내부 깊이를 반영하고, 두 번째 수명 $\ell_{\mathrm{second}}$는 ``경쟁 형성'' 또는 잡음을 반영한다.

\subsection{관절 비율의 의미}

$\mathrm{Artic} = 1 - \ell_{\mathrm{second}}/\ell_{\max}$는 주요 형성이 얼마나 ``우세(dominant)''한지를 측정한다:
\begin{itemize}
    \item $\mathrm{Artic} \approx 1$: 하나의 명확한 주요 형성이 존재하고 경쟁 구조가 없음
    \item $\mathrm{Artic} \approx 0$: 두 개의 비슷한 크기의 구조가 존재 → 형태학적으로 모호
\end{itemize}

\begin{example}[H0 지속 다이어그램의 수동 계산]
\label{ex:h0-persistence}
$5 \times 5$ 격자 위의 형성 $u$를 고려하자. 핵심 값이 다음과 같다고 하자:
\begin{equation}
    u = \begin{pmatrix}
    0.05 & 0.10 & 0.15 & 0.08 & 0.03 \\
    0.12 & 0.45 & 0.72 & 0.38 & 0.06 \\
    0.18 & 0.78 & 0.95 & 0.81 & 0.14 \\
    0.09 & 0.41 & 0.68 & 0.35 & 0.05 \\
    0.04 & 0.08 & 0.12 & 0.07 & 0.02
    \end{pmatrix} \nonumber
\end{equation}

문턱값 $\theta$를 1에서 0으로 내리면서 연결 성분을 추적한다:
\begin{itemize}
    \item $\theta = 0.95$: $\{(3,3)\}$ 탄생 → 성분 A (최대값 지점)
    \item $\theta = 0.81$: $\{(3,4)\}$ 탄생 → 성분 B, 즉시 A와 합병 (인접)
    \item $\theta = 0.78$: $\{(3,2)\}$ 추가 → A에 합류
    \item $\theta = 0.72$: $\{(2,3)\}$ 추가 → A에 합류
    \item $\theta = 0.68$: $\{(4,3)\}$ 추가 → A에 합류
    \item $\theta = 0.45$: $\{(2,2)\}$ 추가 → A에 합류
    \item $\theta = 0.18$: $\{(3,1)\}$ 탄생 → 성분 C, 즉시 A와 합병
    \item 이하 생략...
\end{itemize}

지속 다이어그램: 성분 A의 수명 = $0.95 - 0 = 0.95$ (최종까지 생존). 모든 다른 성분은 탄생 직후 A에 합병되므로 수명이 매우 짧다.

따라서: $\ell_{\max} = 0.95$, $\ell_{\mathrm{second}} \approx 0$ (정확히는 가장 늦게 합병하는 외곽 성분의 짧은 수명).

$c = \bar{u} = 0.302$로, $\mathcal{Q}_{\mathrm{morph}} = \frac{0.95 - 0.302}{1 - 0.302} \times (1 - 0/0.95) \approx 0.93$.

이 높은 $\mathsf{Inside}$ 값은 하나의 명확한 우세 성분(핵심 $(3,3)$ 중심)이 존재함을 반영한다.
\end{example}

\subsection{QM 공리 (QM1--QM4)}

$\mathcal{Q}_{\mathrm{morph}}$는 다음 네 공리를 만족한다:

\begin{theorem}[형태학적 품질 공리]
\label{thm:qm-axioms}
\begin{enumerate}
    \item[\textbf{QM1}] \textbf{단조성}: 형성의 핵심이 강화되면 ($u_{\mathrm{core}}$ 증가) $\mathcal{Q}_{\mathrm{morph}}$ 증가
    \item[\textbf{QM2}] \textbf{문턱화}: 임계 핵심 크기 이하에서 $\mathcal{Q}_{\mathrm{morph}} = 0$
    \item[\textbf{QM3}] \textbf{수축}: 공간적 수축(contraction) 하에서 $\mathcal{Q}_{\mathrm{morph}}$ 비증가
    \item[\textbf{QM4}] \textbf{스케일링}: 그래프 크기에 대해 적절한 스케일링 ($n$-독립적 범위)
\end{enumerate}
\end{theorem}

\begin{proof}[QM1의 증명 스케치]
$u_{\mathrm{core}}$가 증가하면 초수준 집합 여과에서 주요 성분 A의 탄생 시각이 높아진다 ($\theta_{\mathrm{birth}} \uparrow$). 소멸 시각은 변하지 않으므로 ($\theta_{\mathrm{death}} = 0$), $\ell_{\max} = \theta_{\mathrm{birth}} - 0$이 증가한다. $(l_{\max} - c)/(1-c)$에서 $c = \bar{u}$도 증가하나, $l_{\max}$의 증가가 $c$의 증가를 지배한다 (핵심 지점 수 $\ll n$이므로 $c$의 변화는 상대적으로 작다). 따라서 $\mathcal{Q}_{\mathrm{morph}}$ 증가.
\end{proof}


\section{Persist --- 시간적 계승}
\label{sec:persist}

\subsection{두 가지 구현}

$\mathsf{Persist}$는 두 가지 방식으로 계산할 수 있다:

\textbf{(1) 핵심-겹침 근사(Core-Overlap Approximation):}
\begin{equation}
    \mathsf{Persist}_{\mathrm{overlap}}(u_t, u_s) = \frac{\sum_{x} \min(u_t(x),\, u_s(x))}{\max\!\big(\sum_x u_t(x),\, \sum_x u_s(x)\big)}
    \label{eq:persist-overlap}
\end{equation}
이 근사는 전달 핵 계산 없이 두 시각의 응집 장을 직접 비교한다. 부드러운 장 겹침(soft-field overlap)을 측정하며, 계산이 즉각적이다.

\textbf{(2) 전달 기반(Transport-Based):}
\begin{equation}
    \mathsf{Persist}_{\mathrm{tr}}(u_t, u_s) = \frac{\sum_{x \in \mathrm{Core}_t} \sum_{y \in \mathrm{Core}_s} \mathbf{M}_{t \to s}(x,y) \cdot u_s(y)}{\rho_{\mathrm{persist}}}
    \label{eq:persist-transport}
\end{equation}
여기서 $\rho_{\mathrm{persist}}$는 정규화 상수이다. 이 구현은 §\ref{sec:cohesion-fingerprint}--\ref{sec:fixed-point}의 전달 핵 $\mathbf{M}_{t \to s}$를 사용하며, 핵심-핵심 구조적 계승을 정밀하게 측정한다.

\subsection{언제 어떤 구현을 사용하는가}

\begin{center}
\begin{tabular}{lcc}
\hline
\textbf{기준} & \textbf{핵심-겹침} & \textbf{전달 기반} \\
\hline
계산 비용 & $O(n)$ & $O(n^2 \cdot \text{Sinkhorn 반복})$ \\
정밀도 & 정성적 선별 & 엄밀한 구조적 분석 \\
전달 핵 필요 & 아니오 & 예 \\
자기-참조 반영 & 아니오 & 예 \\
사용 상황 & 빠른 탐색 & 이론적 증명, 정밀 진단 \\
\hline
\end{tabular}
\end{center}

실용적 지침: 대부분의 경우 핵심-겹침 근사로 충분하다. 전달 기반 $\mathsf{Persist}$는 T-Persist-1의 조건을 엄밀하게 검증하거나, 약한 체제의 자기-참조 효과를 분석할 때 필요하다.

\begin{example}[두 구현의 비교]
\label{ex:persist-comparison}
$15 \times 15$ 격자에서의 단일 형성, 시각 $t$와 $t+1$ 사이의 지속성:

\textbf{시나리오 A: 안정적 형성 (핵심 불변).}
\begin{itemize}
    \item 핵심-겹침: $\mathsf{Persist}_{\mathrm{overlap}} = 0.94$
    \item 전달 기반: $\mathsf{Persist}_{\mathrm{tr}} = 0.96$
    \item 두 구현의 차이: $0.02$ (미미). 핵심이 거의 이동하지 않으므로 두 방법 모두 높은 지속성을 보고한다.
\end{itemize}

\textbf{시나리오 B: 이동하는 형성 (핵심 2 노드 이동).}
\begin{itemize}
    \item 핵심-겹침: $\mathsf{Persist}_{\mathrm{overlap}} = 0.71$ (겹침 감소)
    \item 전달 기반: $\mathsf{Persist}_{\mathrm{tr}} = 0.89$ (전달이 이동을 추적)
    \item 차이: $0.18$. 전달 기반이 더 정확---핵심-겹침은 이동을 ``질량 소실''로 오판.
\end{itemize}

\textbf{시나리오 C: 약해지는 형성 (핵심 밀도 감소).}
\begin{itemize}
    \item 핵심-겹침: $\mathsf{Persist}_{\mathrm{overlap}} = 0.62$
    \item 전달 기반: $\mathsf{Persist}_{\mathrm{tr}} = 0.58$
    \item 차이: $0.04$. 형성이 제자리에서 약해지는 경우, 두 방법이 비슷한 결과를 준다.
\end{itemize}

결론: 이동이 없는 경우 두 구현의 차이는 미미하다. 이동이 있는 경우 전달 기반이 더 정확하지만, 핵심-겹침도 정성적 선별에는 충분하다.
\end{example}


\section{원시-응집 존재 정리}
\label{sec:proto-cohesion-existence}

\subsection{존재 정리의 조립}

지금까지 별개로 다루어진 네 성분을 조립하면, 다음 존재 정리가 성립한다:

\begin{theorem}[원시-응집 존재]
\label{thm:proto-cohesion-existence}
$\beta/\alpha > 4\lambda_2/|W''(c)|$ (T8-Core 상전이 조건)를 만족하는 매개변수에서, $\Sigma_m$ 위의 에너지 최소점 $u^*$는 네 진단 성분 모두에 대해 비자명한(nontrivial) 하한을 만족한다:
\begin{align}
    \mathsf{Bind}(u^*) &\geq 1 - \sqrt{\mathcal{E}_{\mathrm{cl}}(u^*)/n} > 0 \label{eq:bind-bound} \\
    \mathsf{Sep}(u^*) &= 1 - \mathcal{E}_{\mathrm{sep}}(u^*)/m > 0 \label{eq:sep-bound} \\
    \mathsf{Inside}(u^*) &> 0 \quad \text{($|\mathrm{Core}| \geq 25$일 때)} \label{eq:inside-bound} \\
    \mathsf{Persist}(u^*) &> 0 \quad \text{(T-Persist-1 조건 하)} \label{eq:persist-bound}
\end{align}
\end{theorem}

\begin{proof}[증명의 구조]
\begin{enumerate}
    \item T1(존재): $\Sigma_m$의 컴팩트성 + $\mathcal{E}$의 연속성 → 최소점 $u^*$ 존재
    \item T8-Core(상전이): $\beta/\alpha$ 조건 → $u^*$가 비자명 (균일이 아닌 구조화된 형성)
    \item 식~\eqref{eq:bind-bound}: $\mathsf{Bind}$의 에너지 하한 (Cauchy-Schwarz)
    \item 식~\eqref{eq:sep-bound}: 술어-에너지 다리 (정확한 등식)
    \item 식~\eqref{eq:inside-bound}: 비자명 형성의 H0 지속 ($\ell_{\max} > c$ for structured $u^*$)
    \item 식~\eqref{eq:persist-bound}: T-Persist-1의 5조건 + 2층 집중
\end{enumerate}
\end{proof}

\textbf{해석}: 형성은 ``우연히'' 나타나는 것이 아니다. 상전이 임계점을 넘으면, 에너지 최소화는 \emph{기하학적으로} 네 차원 모두에서 비자명한 형성을 보장한다.

\begin{center}
\fbox{\parbox{0.85\textwidth}{\centering
\textbf{원시-응집의 핵심 의미}\\[6pt]
상전이 이후, 에너지 최소점은 ``자동으로'' 원시-응집을 만족한다.\\
형성의 구조적 품질은 에너지 곡면의 기하학에 내장되어 있다.
}}
\end{center}

% [Figure placeholder]
\begin{figure}[t]
\centering
\fbox{\parbox{0.85\textwidth}{\centering\vspace{5cm}
\textbf{[그림 11.1]} 진단 벡터의 예시.\\
좌: 강한 형성 $\mathbf{d} = (0.94, 0.91, 0.87, 0.93)$ --- 레이더 차트로 네 축 모두 높음.\\
중: 분리가 약한 형성 $\mathbf{d} = (0.88, 0.42, 0.75, 0.90)$ --- Sep 축만 낮음.\\
우: 약한 형성 $\mathbf{d} = (0.35, 0.28, 0.18, 0.55)$ --- 전체적으로 낮음.
\vspace{0.5cm}}}
\caption{세 가지 형성의 진단 벡터. 레이더 차트는 각 형성의 ``프로필''을 직관적으로 보여주며, 약한 차원의 진단을 가능하게 한다.}
\label{fig:diagnostic-examples}
\end{figure}


\section{다섯 가지 검증된 예측}
\label{sec:five-predictions}

SCC 이론은 구체적이고 검증 가능한 예측을 생성한다. 다음 다섯 가지 예측이 수치 실험을 통해 검증되었다.

\subsection{P1: 수축 (Contraction)}

\textbf{예측}: 에너지 최소화 과정에서 형성은 공간적으로 수축한다. 초기 형성의 지지(support)가 최적화 후 줄어들며, 핵심이 강화된다.

\textbf{검증}: 격자 위의 가우시안 범프에서 출발하여 최적화하면, 형성의 유효 반경이 지속적으로 감소하고 핵심 밀도가 증가하는 것이 확인됨.

\textbf{메커니즘}: 이중 우물 에너지 $W(u) = u^2(1-u)^2$이 $u = 0$ 또는 $u = 1$로의 분극을 유도하고, 체적 제약이 총 질량을 고정하므로, 질량이 핵심으로 집중된다.

\subsection{P2: 독립성 (Independence)}

\textbf{예측}: 분리 에너지 계수 $\lambda_{\mathrm{sep}}$을 증가시켜도 $\mathsf{Bind}$가 파괴되지 않는다. 분리 강화가 자기-지지를 저해하지 않는다.

\textbf{검증}: $\lambda_{\mathrm{sep}}$을 0에서 10까지 변화시키는 스윕 실험에서, $\mathsf{Bind}$는 안정적으로 유지됨 (변동 $< 5\%$).

\textbf{메커니즘}: $\mathcal{E}_{\mathrm{cl}}$과 $\mathcal{E}_{\mathrm{sep}}$는 에너지의 독립적 항이며, 최소화 과정에서 $\mathcal{E}_{\mathrm{sep}}$의 기울기는 $\mathcal{E}_{\mathrm{cl}}$의 기울기와 대체로 직교한다.

\subsection{P3: 강화된 체류 (Enhanced Dwell)}

\textbf{예측}: 닫힘 연산자의 존재가 형성의 준안정성을 강화한다. SCC 형성은 순수 Allen-Cahn 에너지의 형성보다 잡음에 더 강하다.

\textbf{검증}: 잡음 주입 실험에서 형성의 ``체류 시간''(형성이 구조를 유지하는 기간)을 비교한 결과, SCC가 순수 Allen-Cahn 대비 \textbf{통계적으로 유의미하게} 더 긴 체류 시간을 보임 ($p = 0.037$).

\textbf{메커니즘}: 닫힘은 형성의 핵심을 자기-강화(self-reinforcing)하여, 헤시안의 최소 고유값을 증가시킨다 (T3/T6-Stability). 이는 에너지 곡면의 극소 분지를 더 깊고 좁게 만든다.

\subsection{P4: 경로 의존성 (Path Dependence)}

\textbf{예측}: 최종 형성의 형태는 초기 조건에 의존한다. 같은 매개변수에서도 다른 초기 조건은 다른 극소점으로 수렴할 수 있다.

\textbf{검증}: 다중 시작(multi-start) 최적화에서 상이한 초기 조건이 상이한 극소점으로 수렴하는 것이 확인됨. 이는 에너지 곡면의 다중 극소점 구조와 일치한다.

\textbf{함의}: 형성은 역사에 의존한다---같은 환경에서도 형성의 ``경험''에 따라 다른 안정 구조가 나타날 수 있다.

\subsection{P5: Sep-before-Inside 순서}

\textbf{예측}: 에너지 최소화 과정에서 $\mathsf{Sep}$이 $\mathsf{Inside}$보다 먼저 문턱을 넘는다. 즉, 외부와의 구별이 내부 구조의 명확화보다 선행한다.

\textbf{검증}: 최적화 궤적을 추적한 결과, $\mathsf{Sep} > 0.5$에 도달하는 반복 횟수가 $\mathsf{Inside} > 0.5$에 도달하는 반복 횟수보다 일관되게 적음이 확인됨.

\textbf{메커니즘}: 구별 연산자 $D_t$는 $u$의 이웃 구조에만 의존하는 국소적 연산이므로, 빠르게 수렴한다. 반면 H0 지속 기반의 $\mathsf{Inside}$는 전역적 위상 구조를 필요로 하므로, 형성이 충분히 안정화된 후에야 높은 값에 도달한다.

\subsection{예측의 정량적 요약}

다섯 예측의 검증 결과를 정량적으로 정리한다:

\begin{center}
\begin{tabular}{lcccc}
\hline
\textbf{예측} & \textbf{측정 지표} & \textbf{예측값} & \textbf{실험값} & \textbf{상태} \\
\hline
P1 수축 & 유효 반경 변화 & $< 0$ & $-23\%$ & 검증 \\
P2 독립성 & Bind 변동 ($\lambda_{\mathrm{sep}}$ 스윕) & $< 10\%$ & $4.2\%$ & 검증 \\
P3 체류 & SCC vs. AC 체류 비 & $> 1$ & $1.38$ ($p=0.037$) & 검증 \\
P4 경로 의존 & 상이한 극소점 수 & $\geq 2$ & 3--5 & 검증 \\
P5 순서 & Sep 문턱 도달 반복 / Inside 문턱 도달 반복 & $< 1$ & $0.42$ & 검증 \\
\hline
\end{tabular}
\end{center}

\subsection{예측과 이론의 관계}

다섯 예측은 단순히 ``확인된 사실''이 아니라, 이론의 구조적 특성에서 연역된 것이다:

\begin{itemize}
    \item \textbf{P1 (수축)}: 에너지 범함수의 이중 우물 + 체적 제약의 직접적 귀결
    \item \textbf{P2 (독립성)}: 네 에너지 항의 독립성 원리의 검증
    \item \textbf{P3 (강화된 체류)}: 비-멱등 닫힘의 안정성 이점 (T3/T6-Stability)
    \item \textbf{P4 (경로 의존)}: 에너지 곡면의 다중 극소점 구조
    \item \textbf{P5 (Sep-before-Inside)}: 국소 연산(구별) vs. 전역 연산(H0 지속)의 수렴 속도 차이
\end{itemize}

이 중 P3는 특히 중요하다: SCC 이론이 순수 Allen-Cahn(또는 표준 위상장 모형)과 \emph{구별되는} 예측을 생성하고, 이것이 통계적으로 유의미하게 검증되었기 때문이다. $p = 0.037 < 0.05$는 ``SCC의 닫힘이 준안정성을 강화한다''는 주장이 우연이 아님을 보여준다.

% [Figure placeholder]
\begin{figure}[t]
\centering
\fbox{\parbox{0.85\textwidth}{\centering\vspace{5cm}
\textbf{[그림 11.2]} $\mathcal{Q}_{\mathrm{morph}}$의 H0 지속 다이어그램.\\
좌: 강한 형성의 지속 다이어그램---하나의 긴 막대가 우세.\\
우: 약한 형성---여러 막대가 비슷한 길이, $\mathrm{Artic} \approx 0$.
\vspace{0.5cm}}}
\caption{H0 지속 다이어그램에 의한 형태학적 품질 평가. 강한 형성은 하나의 우세한 성분을 가지며, 약한 형성은 경쟁 구조가 존재한다.}
\label{fig:qmorph-persistence}
\end{figure}

% [Figure placeholder]
\begin{figure}[t]
\centering
\fbox{\parbox{0.85\textwidth}{\centering\vspace{5cm}
\textbf{[그림 11.3]} 다섯 가지 예측의 검증.\\
(a) P1: 수축---유효 반경 감소. (b) P2: $\lambda_{\mathrm{sep}}$ 스윕에서 Bind 안정.\\
(c) P3: SCC vs. Allen-Cahn 체류 시간 비교. (d) P4: 다중 시작의 상이한 수렴.\\
(e) P5: Sep과 Inside의 수렴 속도 비교.
\vspace{0.5cm}}}
\caption{SCC 이론의 다섯 가지 예측(P1--P5) 검증 결과. 모든 예측이 수치 실험에 의해 확인되었다.}
\label{fig:prediction-verification}
\end{figure}


\section{요약}
\label{sec:diagnostic-summary}

이 장에서 확립된 핵심 사항:

\begin{itemize}
    \item \textbf{진단 벡터}: $\mathbf{d} = (\mathsf{Bind}, \mathsf{Sep}, \mathsf{Inside}, \mathsf{Persist}) \in [0,1]^4$는 형성의 구조적 품질을 네 독립적 차원에서 연속적으로 평가한다.

    \item \textbf{Bind}: $\ell^2$ 닫힘 잔차 기반. T-Bind-Proj(사영 단조성), T-Bind-Full(에너지 최소점 하한).

    \item \textbf{Sep}: $u$-가중 구별 평균. 정확한 등식 $\mathsf{Sep} = 1 - \mathcal{E}_{\mathrm{sep}}/m$ (술어-에너지 다리).

    \item \textbf{Inside}: H0 지속 기반 $\mathcal{Q}_{\mathrm{morph}}$. QM1--4 공리 만족.

    \item \textbf{Persist}: 핵심-겹침(빠른 탐색) 또는 전달 기반(정밀 분석)의 이중 구현.

    \item \textbf{존재 정리}: 상전이 이후, 에너지 최소점은 자동으로 네 차원 모두에서 비자명한 원시-응집을 만족한다.

    \item \textbf{다섯 예측}: P1(수축), P2(독립성), P3(강화된 체류, $p=0.037$), P4(경로 의존성), P5(Sep-before-Inside) 모두 검증됨.
\end{itemize}

제III부를 마무리하며, 우리는 SCC 이론의 형식적 구조를 완성했다: 단일 형성(제II부) → 다중 형성(제9장) → 시간적 전달(제10장) → 통합 진단(이 장). 제IV부에서는 이 이론을 구현하고, 인지과학적 연결을 탐구하며, 이론의 한계와 미래 방향을 논의한다.


\section*{제III부 회고: 형식에서 현상으로}

제III부는 이론의 가장 도전적인 부분을 다루었다. 핵심 교훈을 세 가지로 정리한다.

\textbf{첫째, 다중 형성은 에너지의 문제가 아니라 역학의 문제이다.} 전역 에너지 최소점은 항상 $K = 1$이다. 복수의 형성이 공존하는 것은 운동학적 장벽에 의해 보호되는 준안정 현상이다. 이 관점의 전환---열역학에서 운동학으로---은 이론의 가장 중요한 개념적 혁신 중 하나이다.

\textbf{둘째, 시간적 정체성은 라벨이 아니라 구조적 계승이다.} 추적 ID 없이 정체성을 정의하는 것은 가능하며, 응집 지문과 자기-참조적 전달 핵이 그 수학적 실현이다. 이 접근은 분열, 합병, 점진적 변형 등 전통적 추적의 난제를 자연스럽게 처리한다.

\textbf{셋째, 진단 벡터는 이론의 ``인터페이스''이다.} 네 성분의 연속적 평가는 형성의 품질을 다차원적으로 진단할 수 있게 하며, 다섯 가지 검증된 예측은 이론의 경험적 기반을 제공한다. 특히 P3(강화된 체류, $p = 0.037$)는 SCC가 기존 이론과 구별되는 고유한 예측을 생성함을 보여준다.

제IV부에서는 이 형식적 구조를 파이썬 코드로 구현하고(제12장), 수치 실험으로 검증하며(제13장), 인지과학과의 연결을 탐구하고(제14장), 열린 문제와 미래 방향을 논의한다(제15장).
